\section{Ethical Use of Data}
\label{sec:ethics}
We use only public data, the verifiability judgments of Liu et al.~\cite{Liu_2023} (MIT licence); no human participants were involved, and the 2026 spot checks were labelled by AI agents on locally stored text. Requests to cited sites: (i) our identified client, once per URL for the measurement (at least five seconds between requests to a host), and once per run for the pages of the \eThreeN E3 answers (three full runs; explainer runs on one and twice on five answers), the \eFourN E4 answers (four runs) and example answers citing six sites (about twenty test and video runs, plus a few pilot and connectivity tests), two seconds apart per host (a host's robots.txt and first page request were unspaced until a late fix); sites that ask automated clients not to fetch a page were never requested (redirect targets were not re-checked); (ii) one visit with an ordinary browser under automation for \browserN refused, failing or empty pages (\browserTwice twice; \browserUnused dead training-split pages opened under a broader rule are not used), without solving challenges, logging in or bypassing paywalls; pages reached only after a detected bot interstitial count as bot walls; (iii) one urlhealth check per fetchable development/test URL, run as published (HEAD, a GET after some errors, an archive lookup after a 404; browser user agent), rate-limited like all our archive lookups. Page bodies stay local: the public explorer shows the dataset's evidence and at most one current-page sentence (up to 300 characters) per pair. Check mode processes the answer locally, stores nothing, and requests only the cited pages, their robots.txt files and archive records (and, on request, one browser visit; one such test was run).
